\documentclass[10pt,a4paper]{amsart}
\usepackage[margin=19mm]{geometry}
\usepackage{amsmath,amssymb,mathtools}
\usepackage[scr=rsfs,cal=euler]{mathalfa}
\usepackage{xfrac}
\usepackage[unicode,hidelinks,bookmarks=false]{hyperref}
\newcommand{\ie}{i.e.\ }
\newcommand{\cf}{cf.\ }
\newcommand{\resp}{respectively\ }
\newcommand{\Subsection}{\subsection}
\providecommand{\nobreakdash}{\nobreak\hbox{-}}
\setlength{\emergencystretch}{2em}
\multlinegap0pt
\usepackage{amssymb}
\usepackage{bm}
\usepackage[shortlabels]{enumitem}
\usepackage{bbm}
\usepackage{mathtools}
\usepackage{xcolor}
\usepackage{tikz}
\usepackage{multirow}
\usepackage{float}
\newtheorem{lemma}{Lemma}
\DeclareMathOperator{\Hom}{Hom}
\DeclareMathOperator{\Ker}{Ker}
\title{The Global Jacquet--Langlands Correspondence via Tensor Products}
\author{Jun Yang}
\date{Source: arXiv:2602.08053v1 (CC BY 4.0)}
\begin{document}
\maketitle

\noindent\emph{Source and licence.} The Global Jacquet--Langlands Correspondence via Tensor Products. Version arXiv:2602.08053v1 (CC BY 4.0), distributed by arXiv under the Creative Commons Attribution 4.0 International licence, http://creativecommons.org/licenses/by/4.0/, as recorded in the arXiv licence metadata for this exact version and shown on the abstract page https://arxiv.org/abs/2602.08053v1. Full licence notice: \texttt{licenses/arxiv-2602.08053-CC-BY-4.0.txt}.\par\smallskip

\noindent\textit{Editorial scope.} Counted unit: the article's lemma ltenHsim (source0.tex lines 617--625). The article's complete original proof of it is delivered with it (source0.tex lines 627--651, one \texttt{proof} environment of 25 lines). No mathematics is written, repaired or re-derived: the statement and the proof are the article's own lines, in the article's own notation, and no retained line is substituted or re-flowed.  No further source-local context is needed: every cross-reference of every retained line resolves inside the two delivered spans.  Nothing was omitted from any retained span, so the package carries no omission record.  No retained line contains a citation, so the wrapper carries no bibliography and no bibliography entry.  No retained line contains a figure environment, an image-inclusion command or drawing code, so no image is delivered and the package declares no asset dependency.  Editorial formatting only, in addition: an AMS \texttt{amsart} preamble that restates the article's own declarations and macros used by the retained lines, namely the environments \texttt{lemma} on separate counters, together with the standard packages those lines need.  The retained environments print in this document as Lemma 1 in order of appearance, whereas the article prints the same items with the numbering of its own full document.  The complete native source of the article as distributed in the arXiv e-print for this exact version is bundled unchanged as \texttt{sources/194227/source0.tex}.
\begin{lemma}\label{ltenHsim}
Let $(G,H)$ be a similitude dual pair over a local field $F$, and let $\Omega$ be the induced representation of $G\times H$.
For an irreducible representation $\pi$ of $G$, we have
\[
\Theta(\pi)\cong\pi^{\vee}\otimes_{\mathcal{H}(G)}\Omega
\quad \text{(resp.\ } \Theta(\pi)\cong\pi^{\vee}\otimes_{\mathcal{H}(\mathfrak{g},K)}\Omega\text{)}
\]
as representations of $H$, according to whether $F$ is non-archimedean (resp.\ archimedean).
\end{lemma}

\begin{proof}
We write $\mathcal{H}$ for $\mathcal{H}(G)$ (resp.\ $\mathcal{H}(\mathfrak{g},K)$).
Let $\mathcal{T}$ be the underlying space of $\Omega$.
Recall that
\[
\mathcal{T}(\pi)=\bigcap_{T\in \Hom_{G}(\mathcal{T},V_{\pi})}\Ker(T).
\]
Consider the short exact sequence
\[
0\to \mathcal{T}(\pi)\to \mathcal{T} \to V_{\pi}\otimes_{\mathbb{C}} V_{\Theta(\pi)}\to 0.
\]
Tensoring $V_{\pi}^{\vee}$ over $\mathcal{H}$ yields the exact sequence
\[
V_{\pi}^{\vee}\otimes_{\mathcal{H}}\mathcal{T}(\pi)\to V_{\pi}^{\vee}\otimes_{\mathcal{H}}\mathcal{T} \to V_{\pi}^{\vee}\otimes_{\mathcal{H}}V_{\pi}\otimes_{\mathbb{C}} V_{\Theta(\pi)}\to 0.
\]
Since $V_{\pi}^{\vee}\otimes_{\mathcal{H}}V_{\pi}\cong \Hom_{G}(V_{\pi},V_{\pi}^{\vee})\cong\mathbb{C}$, we obtain
\[
V_{\pi}^{\vee}\otimes_{\mathcal{H}}\mathcal{T}(\pi)\to V_{\pi}^{\vee}\otimes_{\mathcal{H}}\mathcal{T} \to  V_{\Theta(\pi)}\to 0.
\]
If $V_{\pi}^{\vee}\otimes_{\mathcal{H}}\mathcal{T}(\pi)\neq 0$, then
$V_{\pi}\subset \mathcal{T}(\pi)$ since
$V_{\pi}^{\vee}\otimes_{\mathcal{H}}\mathcal{T}(\pi)\subset \Hom_{\mathcal{H}}(V_{\pi},\mathcal{T}(\pi))$,
which contradicts the definition of $\mathcal{T}(\pi)$.
Hence $\Theta(\pi)\cong\pi^{\vee}\otimes_{\mathcal{H}}\Omega$.
\end{proof}

\end{document}
